\documentclass[conference]{IEEEtran}
\IEEEoverridecommandlockouts

\usepackage{amsmath,amssymb}
\usepackage{cite}
\usepackage{booktabs}
\usepackage{array}
\usepackage{url}

\begin{document}

\title{Baseline-Referenced Feature Extraction for Task-Evoked Changes in Single-Channel EEG}

\author{
\IEEEauthorblockN{Author Name}
\IEEEauthorblockA{Affiliation\\
Email}
}

\maketitle

\begin{abstract}
This paper presents a baseline-referenced method for extracting task-evoked changes from single-channel frontal electroencephalography (EEG). Raw EEG is segmented into two-second windows and transformed into spectral power, relative power, peak, entropy, and cross-band-ratio features using multitaper power spectral density estimation. After signal-quality screening and electromyographic artifact control, adjacent windows are aggregated into temporal blocks. For each feature, the principal task descriptor is the difference between the task mean and the participant-specific eyes-closed baseline mean. This feature delta is supplemented by percent change, standardized change, Cohen effect size, directional consistency, and inferential uncertainty. Welch tests, false-discovery-rate adjustment, dependence-aware Fisher aggregation, and permutation testing quantify the reliability of the observed baseline-to-task changes. The resulting representation emphasizes within-participant task modulation rather than absolute EEG magnitude, thereby reducing sensitivity to stable inter-individual and device-dependent differences.
\end{abstract}

\begin{IEEEkeywords}
EEG, single-channel EEG, baseline normalization, feature extraction, task-evoked change, multitaper spectral analysis, effect size, cognitive-state analysis.
\end{IEEEkeywords}

\section{Introduction}
Consumer-grade single-channel EEG provides a practical route for repeated cognitive-state measurement, but absolute spectral values vary substantially across participants, recording sessions, electrode contact conditions, and hardware gain. Consequently, an isolated task measurement is difficult to interpret. A participant-specific baseline provides a physiological reference against which task-related modulation can be expressed.

The central quantity in this work is therefore not the absolute feature value, but its change from baseline. Let $\mu_f^{(B)}$ denote the mean baseline value of feature $f$ and $\mu_f^{(t)}$ the corresponding mean during task $t$. The primary extracted task feature is
\begin{equation}
\Delta_f^{(t)}=\mu_f^{(t)}-\mu_f^{(B)}.
\end{equation}
The vector $\boldsymbol{\Delta}^{(t)}$ summarizes the participant's task-evoked spectral response. Its sign indicates the direction of modulation, while its magnitude is characterized using absolute, relative, and standardized measures. Statistical procedures are used to assess whether the observed changes are reliable given temporal variability and correlation among EEG features.

\section{Method Overview}
The scientific processing sequence is:
\begin{enumerate}
\item acquire single-channel frontal EEG at 512 Hz;
\item segment baseline and task recordings into non-overlapping two-second windows;
\item estimate spectral density and extract window-level features;
\item reject unusable baseline windows and suppress artifact-sensitive gamma features;
\item aggregate consecutive feature windows into temporal blocks;
\item compute baseline and task feature means;
\item derive baseline-referenced feature deltas and normalized change measures;
\item quantify the reliability and direction of each change;
\item compare delta profiles across cognitive tasks.
\end{enumerate}

\section{Signal Acquisition and Filtering}
Single-channel frontal EEG is sampled at $f_s=512$ Hz. For real-time visualization, raw samples are processed by a cascaded second-order IIR chain consisting of a 1 Hz high-pass filter, a 45 Hz low-pass filter, and a 50 Hz notch filter. Each section follows the direct-form I biquad:
\begin{equation}
y[n]=b_0x[n]+b_1x[n-1]+b_2x[n-2]-a_1y[n-1]-a_2y[n-2].
\end{equation}
The high-pass and low-pass filters use Butterworth biquad coefficients obtained through the bilinear transform. For cutoff $f_c$ and sampling frequency $f_s$, let
\begin{equation}
k=\tan\left(\frac{\pi f_c}{f_s}\right),\quad
g=\frac{1}{1+\sqrt{2}k+k^2}.
\end{equation}
The high-pass section uses
\begin{equation}
b_0=g,\quad b_1=-2g,\quad b_2=g,
\end{equation}
\begin{equation}
a_1=2(k^2-1)g,\quad a_2=(1-\sqrt{2}k+k^2)g.
\end{equation}
The low-pass section uses
\begin{equation}
b_0=k^2g,\quad b_1=2k^2g,\quad b_2=k^2g,
\end{equation}
with the same $a_1$ and $a_2$ form. The notch filter is centered at $f_0=50$ Hz with quality factor $Q=35$:
\begin{equation}
\omega_0=\frac{2\pi f_0}{f_s},\quad
\alpha=\frac{\sin \omega_0}{2Q},\quad a_0=1+\alpha .
\end{equation}
The normalized notch coefficients are
\begin{equation}
b_0=\frac{1}{a_0},\quad b_1=\frac{-2\cos \omega_0}{a_0},\quad b_2=\frac{1}{a_0},
\end{equation}
\begin{equation}
a_1=\frac{-2\cos \omega_0}{a_0},\quad a_2=\frac{1-\alpha}{a_0}.
\end{equation}
Feature extraction is performed from the recorded raw samples using the spectral procedure described below.

\section{Windowing and Spectral Estimation}
\subsection{Segmentation}
The raw sequence is partitioned into non-overlapping windows:
\begin{equation}
\mathbf{x}_w=\{x[wN],x[wN+1],\ldots,x[wN+N-1]\},
\end{equation}
where $N=1024$ and $f_s=512$ Hz. Each window therefore spans 2 s. The window is demeaned:
\begin{equation}
\tilde{x}_w[n]=x_w[n]-\frac{1}{N}\sum_{m=0}^{N-1}x_w[m].
\end{equation}

\subsection{DPSS Multitaper PSD}
A multitaper PSD is computed with $K=3$ discrete prolate spheroidal sequence tapers and time-bandwidth product $NW=2.5$. For taper $v_k[n]$, the tapered Fourier transform is
\begin{equation}
X_k[f]=\sum_{n=0}^{N-1}\tilde{x}_w[n]v_k[n]e^{-j2\pi fn/N}.
\end{equation}
The PSD estimate is
\begin{equation}
\hat{S}[f]=\frac{1}{K}\sum_{k=1}^{K}\frac{|X_k[f]|^2}{f_sN+\epsilon},
\end{equation}
where $\epsilon$ prevents numerical division errors. A Hann-tapered periodogram provides a fallback estimate when multitaper estimation is unavailable.

\subsection{Noise-Floor Subtraction}
For each PSD, a robust noise floor is estimated as the 10th percentile:
\begin{equation}
\eta_w = Q_{0.10}\left(\hat{S}_w[f]\right).
\end{equation}
The SNR-adjusted PSD is
\begin{equation}
\hat{S}^{+}_w[f]=\max(\hat{S}_w[f]-\eta_w,0).
\end{equation}
This reduces the influence of broadband floor shifts before band integration.

\section{Quality Control and Artifact Handling}
\subsection{Baseline Window Rejection}
Baseline windows are screened before the baseline reference distribution is estimated. A window is rejected as a flatline when the robust scale is too small:
\begin{equation}
\sigma_{\mathrm{MAD}}=1.4826\cdot\mathrm{median}\{|x[n]-\mathrm{median}(x)|\}<0.5.
\end{equation}
It rejects likely not-worn windows using spectral criteria. Let
\begin{equation}
R_{\mathrm{neural}}=\frac{\sum_{0.5\le f\le 13}\hat{S}[f]}{\sum_f\hat{S}[f]},
\quad
R_{\mathrm{HF}}=\frac{\sum_{f\ge 30}\hat{S}[f]}{\sum_f\hat{S}[f]}.
\end{equation}
The log-log spectral slope $m$ is estimated over 1--40 Hz. A window is rejected when low-frequency neural content is insufficient, high-frequency power is excessive, or the spectrum is implausibly flat. Extreme amplitude outliers also trigger artifact rejection when more than 5 percent of samples exceed 20 robust scales.

\subsection{EMG Guard for Gamma Features}
Frontal single-channel EEG is vulnerable to facial EMG contamination, especially in beta/gamma ranges \cite{whitham2007,goncharova2003}. The artifact indicator is
\begin{equation}
R_{\mathrm{EMG}}=\frac{P_{35-45}}{P_{20-30}+\epsilon},
\end{equation}
and a high-frequency log-log slope over 20--45 Hz:
\begin{equation}
\log \hat{S}[f] = m_{\mathrm{HF}}\log f + b .
\end{equation}
The EMG guard fires when
\begin{equation}
R_{\mathrm{EMG}}>1.2 \quad \mathrm{or} \quad m_{\mathrm{HF}}>-0.6 .
\end{equation}
If triggered, gamma features are set to zero and marked as unevaluated, preventing high-frequency muscle activity from being interpreted as neural gamma engagement.

\section{Feature Extraction}
\subsection{Frequency Bands}
Spectral features are computed from the following bands:
\begin{table}[h]
\centering
\caption{Raw EEG Feature Bands}
\begin{tabular}{lcc}
\toprule
Band & Lower Hz & Upper Hz\\
\midrule
Delta & 0.5 & 4\\
Theta & 4 & 8\\
Alpha & 8 & 13\\
Beta & 13 & 30\\
Gamma & 30 & 45\\
Theta1 & 4 & 6\\
Theta2 & 6 & 8\\
Beta1 & 13 & 20\\
Beta2 & 20 & 30\\
\bottomrule
\end{tabular}
\end{table}

When precomputed band-power observations are used, the available delta, theta, low-alpha, high-alpha, low-beta, high-beta, low-gamma, and mid-gamma components are preserved.

\subsection{Band Power and Relative Power}
For band $B=[f_1,f_2]$, raw and SNR-adjusted powers are
\begin{equation}
P_B^{\mathrm{raw}}=\int_{f_1}^{f_2}\hat{S}[f]\,df,\quad
P_B=\int_{f_1}^{f_2}\hat{S}^{+}[f]\,df .
\end{equation}
Band power is estimated by trapezoidal numerical integration. Relative power is
\begin{equation}
P_B^{\mathrm{rel}}=\frac{P_B}{\sum_{B'}P_{B'}+\epsilon},
\end{equation}
or, when SNR power is absent, by normalizing against the window variance.

\subsection{Peak and Entropy Descriptors}
The normalized spectrum is
\begin{equation}
Z[f]=\frac{\hat{S}[f]-\eta}{\eta+\epsilon}.
\end{equation}
For each band, peak frequency and peak amplitude are
\begin{equation}
f_B^\ast=\arg\max_{f\in B} Z[f],\quad A_B^\ast=\max_{f\in B} Z[f].
\end{equation}
Peak relative amplitude is computed as the peak amplitude divided by the mean band-normalized spectrum. Spectral entropy is
\begin{equation}
H_B=-\sum_{f\in B}p_f\log_2p_f,\quad
p_f=\frac{Z[f]}{\sum_{u\in B}Z[u]+\epsilon}.
\end{equation}

\subsection{Cross-Band Ratios}
Four cross-band ratios are derived:
\begin{equation}
\rho_{\alpha/\theta}=\frac{P_\alpha}{P_\theta+\epsilon},\quad
\rho_{\beta/\alpha}=\frac{P_\beta}{P_\alpha+\epsilon},
\end{equation}
\begin{equation}
\rho_{\beta2/\beta1}=\frac{P_{\beta2}}{P_{\beta1}+\epsilon},\quad
\rho_{\theta2/\theta1}=\frac{P_{\theta2}}{P_{\theta1}+\epsilon}.
\end{equation}
These ratios are used as interpretable markers of cognitive load, attention, and task engagement.

\section{Baseline Reference and Block Aggregation}
\subsection{Eyes-Closed Baseline}
Eyes-closed EEG is used as the participant-specific reference condition because it provides a stable resting comparison for alpha- and theta-related task modulation \cite{klimesch1999}. Eyes-open observations may be retained for descriptive comparison but are not included in the primary baseline estimate.

\subsection{Effective Blocks}
To reduce temporal dependence among adjacent 2 s windows, windows are averaged into non-overlapping blocks. The default block duration is 8 s, so each block contains four feature windows:
\begin{equation}
\bar{z}_{b,f}=\frac{1}{M}\sum_{i=1}^{M}z_{bM+i,f},\quad M=4.
\end{equation}
All Welch tests, permutations, and task summaries are computed on these block-level feature means.

\subsection{Unequal Block Counts}
Task and baseline blocks are not downsampled or randomly equalized. If the task produces $n_a$ blocks and the baseline produces $n_b$ blocks, all $n_a+n_b$ available blocks are retained. Unequal-sample statistical procedures preserve these original group sizes and avoid discarding valid observations when recording durations differ.

\section{Baseline-Referenced Feature Delta}
\subsection{Primary Delta Representation}
For feature $f$, let $\{a_i\}_{i=1}^{n_a}$ denote task blocks and $\{b_j\}_{j=1}^{n_b}$ denote baseline blocks. Their means are
\begin{equation}
\mu_f^{(t)}=\frac{1}{n_a}\sum_{i=1}^{n_a}a_i,\quad
\mu_f^{(B)}=\frac{1}{n_b}\sum_{j=1}^{n_b}b_j.
\end{equation}
The primary extracted task feature is the baseline-referenced delta
\begin{equation}
\Delta_f^{(t)}=\mu_f^{(t)}-\mu_f^{(B)}.
\end{equation}
A positive delta indicates an increase relative to baseline, whereas a negative delta indicates suppression. For task $t$, the complete response is represented by
\begin{equation}
\boldsymbol{\Delta}^{(t)}
=\left[\Delta_1^{(t)},\Delta_2^{(t)},\ldots,\Delta_m^{(t)}\right]^\mathsf{T}.
\end{equation}
This within-participant differencing is intended to reduce the influence of stable absolute-power differences and emphasize task-related modulation.

\subsection{Normalized Delta Measures}
Because EEG features differ in scale and units, the absolute delta is accompanied by percent change, baseline-standardized change, and log-ratio change:
\begin{equation}
\Delta_{f,\%}^{(t)}
=100\frac{\mu_f^{(t)}-\mu_f^{(B)}}{|\mu_f^{(B)}|+\epsilon},
\end{equation}
\begin{equation}
z_f^{(t)}=\frac{\mu_f^{(t)}-\mu_f^{(B)}}{s_f^{(B)}+\epsilon},
\quad
L_f^{(t)}=\log_2\left(\frac{|\mu_f^{(t)}|}{|\mu_f^{(B)}|+\epsilon}+\epsilon\right).
\end{equation}
The absolute delta preserves the original feature units, percent change supports within-feature interpretation, and standardized measures support comparison across heterogeneous feature families.

\subsection{Effect Size}
Cohen's effect size characterizes the magnitude of the baseline-to-task shift relative to block variability \cite{cohen1988}:
\begin{equation}
d_f=\frac{\mu_f^{(t)}-\mu_f^{(B)}}{\sqrt{(s_t^2+s_B^2)/2}+\epsilon}.
\end{equation}
Feature-dependent practical-change thresholds are 0.25 for alpha effect sizes, 0.35 for beta effect sizes, and 0.30 otherwise. Percent-change thresholds are 5 percent for relative-power or ratio features and 10 percent for absolute-power features.

\section{Statistical Reliability of Feature Deltas}
\subsection{Welch's t-Test}
The reliability of $\Delta_f^{(t)}$ is assessed using Welch's statistic:
\begin{equation}
t_f=\frac{\bar{a}_f-\bar{b}_f}{\sqrt{s_a^2/n_a+s_b^2/n_b}}.
\end{equation}
The Welch-Satterthwaite degrees of freedom are
\begin{equation}
\nu_f=\frac{(s_a^2/n_a+s_b^2/n_b)^2}
{\frac{(s_a^2/n_a)^2}{n_a-1}+\frac{(s_b^2/n_b)^2}{n_b-1}}.
\end{equation}
At least two observations are required in each group. A group is treated as near-constant when its range is no greater than
\begin{equation}
\tau=\max\left(10^{-12},10^{-9}\max(1,|x_{\min}|,|x_{\max}|,|\bar{x}|)\right).
\end{equation}
If either group is near-constant, the feature is assigned $t_f=0$ and $p_f=1$, preventing unstable or artificial significance.

\subsection{Directional Priors}
Task-specific directional expectations are applied after the two-sided Welch test. For example, mental arithmetic expects alpha decrease and beta/gamma or beta-alpha-ratio increase; visual imagery expects alpha and alpha-theta-ratio increase; working memory expects theta increase, alpha decrease, and beta-ratio increase. If the expected direction is ``up'', then
\begin{equation}
p_f^{(1)}=
\begin{cases}
p_f^{(2)}/2, & t_f>0,\\
1-p_f^{(2)}/2, & t_f\le 0,
\end{cases}
\end{equation}
and analogously for ``down'' with $t_f<0$. A feature is considered directionally valid only when the observed sign matches the expectation.

\section{Multiple Comparisons and Feature Gating}
\subsection{Benjamini-Hochberg FDR}
Benjamini-Hochberg $q$ values are computed to characterize multiplicity across features \cite{benjamini1995}. Given sorted $p$ values $p_{(1)}\le\cdots\le p_{(m)}$, adjusted values are obtained by the monotone rule
\begin{equation}
q_{(i)}=\min_{j\ge i}\frac{m}{j}p_{(j)}.
\end{equation}

\subsection{Correlation Guard}
Because EEG features are correlated, a Spearman-rank feature correlation matrix is estimated from the retained baseline blocks. An effective feature count is then estimated using eigenvalues:
\begin{equation}
m_{\mathrm{eff}}=\sum_{i=1}^{m}\min(\max(\lambda_i,0),1).
\end{equation}
The local alpha is shrunk by
\begin{equation}
g=\mathrm{clip}\left(\frac{m_{\mathrm{eff}}}{m},0.05,1.0\right),
\quad
\alpha_{\mathrm{local}}=g\alpha .
\end{equation}
This guard makes feature gating more conservative when many features are redundant.

\subsection{Decision Rule}
For each feature, significance is reported if the direction is valid and at least one of the following passes: one-sided $p$ threshold, Cohen effect-size threshold, or percent-change threshold. The BH flag is recorded separately for transparency.

\section{Global Task Evidence}
\subsection{Fisher Combination with Kost-McDermott Correction}
Feature-level $p$ values are combined with Fisher's statistic \cite{fisher1925}:
\begin{equation}
X=-2\sum_{f=1}^{m}\log p_f.
\end{equation}
Because spectral features are dependent, the Kost-McDermott correction is applied \cite{kost2002}. Pairwise feature correlations $r_{ij}$ are converted into covariance contributions:
\begin{equation}
C_{ij}=3.263r_{ij}+0.710r_{ij}^2+0.027r_{ij}^3 .
\end{equation}
The corrected variance is
\begin{equation}
\sigma^2=4m+2\sum_{i<j}C_{ij},\quad \mu=2m.
\end{equation}
The scale and degrees of freedom are
\begin{equation}
c=\frac{\sigma^2}{2\mu},\quad \nu_{\mathrm{KM}}=\frac{2\mu^2}{\sigma^2}.
\end{equation}
The adjusted statistic is $X_{\mathrm{KM}}=X/c$ and its $p$ value is evaluated from the upper tail of a chi-square distribution with $\nu_{\mathrm{KM}}$ degrees of freedom.

\subsection{SumP Permutation Test}
The SumP statistic is
\begin{equation}
S_{\mathrm{obs}}=\sum_{f=1}^{m}p_f .
\end{equation}
Task and baseline blocks are pooled and labels are permuted with fixed seed 42 for $B=1000$ permutations. Each permutation restores the original task size $n_a$ and baseline size $n_b$; no equal-count subsampling is performed. Near-constant features contribute $p=1$ rather than an unstable test result. The empirical left-tail probability is
\begin{equation}
p_{\mathrm{perm}}=\frac{1+\sum_{b=1}^{B}\mathbf{1}\{S_b\le S_{\mathrm{obs}}\}}{B+1}.
\end{equation}
This test provides a distribution-free complement to the corrected Fisher statistic.

\subsection{Composite and Cosine Scores}
The composite evidence score is
\begin{equation}
C=\sum_{f=1}^{m}-\log_{10}(\max(q_f,10^{-12})).
\end{equation}
Cosine similarity is also computed between baseline and task mean feature vectors:
\begin{equation}
\cos(\theta)=\frac{\mathbf{b}\cdot\mathbf{a}}
{\|\mathbf{b}\|_2\|\mathbf{a}\|_2+\epsilon},
\quad D_{\cos}=1-\cos(\theta).
\end{equation}
A permutation procedure shuffles the baseline vector to estimate a cosine-distance $p$ value.

\section{Across-Task Delta Comparison}
When multiple tasks are recorded, each feature is ranked across tasks using the absolute task-versus-baseline effect size. This identifies which cognitive condition produces the largest modulation relative to the same participant-specific baseline. When at least two tasks contain data, the Kruskal-Wallis test compares per-window feature values grouped by task \cite{kruskal1952}:
\begin{equation}
H=\frac{12}{N(N+1)}\sum_{g=1}^{G}n_g(\bar{R}_g-\bar{R})^2 ,
\end{equation}
followed by BH-FDR adjustment over feature-level omnibus $p$ values. Cross-task task-level Fisher/KM $p$ values are additionally adjusted with Holm-Bonferroni \cite{holm1979}:
\begin{equation}
p_{(i)}^{\mathrm{Holm}}=\max_{j\le i}(m-j+1)p_{(j)}.
\end{equation}
The omnibus test evaluates whether at least one task condition differs from the others, while the baseline-referenced delta and effect-size rankings retain the direction and magnitude of each task response.

\section{Scientific Interpretation}
The resulting feature-delta profile is interpreted using the following evidence:
\begin{itemize}
\item the sign and magnitude of $\Delta_f^{(t)}$;
\item percent and standardized change relative to baseline;
\item block-level variability and Cohen effect size;
\item directional consistency with task-specific hypotheses;
\item feature-level and global statistical reliability;
\item baseline signal quality and gamma-artifact sensitivity.
\end{itemize}

The delta profile supports statements about task-related modulation of measured EEG features. It does not directly establish stable psychological traits, occupational abilities, source-localized neural activity, or clinical diagnoses.

\section{Discussion}
The principal contribution of this method is the construction of a task-specific feature representation from changes relative to a personal baseline. Absolute EEG power is affected by stable participant characteristics, sensor placement, impedance, and acquisition gain. Baseline differencing reduces the contribution of these stable offsets and shifts interpretation toward within-session task modulation. The extracted delta vector preserves feature-specific direction, while percent change, standardized change, and Cohen effect size provide complementary measures of magnitude.

Multitaper PSD reduces variance and leakage in short windows \cite{thomson1982,percival1993}; EMG-aware gamma suppression reflects known contamination of scalp EEG by muscle activity \cite{whitham2007,goncharova2003}; Welch testing handles unequal variance and unequal block counts \cite{welch1947}; BH-FDR characterizes multiplicity across features \cite{benjamini1995}; Fisher aggregation summarizes distributed evidence \cite{fisher1925}; and Kost-McDermott correction addresses correlated feature evidence \cite{kost2002}.

The main limitation is spatial specificity. A single frontal/frontopolar-compatible channel cannot support connectivity, source localization, lateralized motor claims, or clinical diagnosis. Baseline differencing also does not remove transient artifacts or time-varying non-neural effects. Baseline raw windows therefore receive explicit signal-quality rejection, while gamma suppression and degenerate-series handling reduce selected risks during task analysis.

\section{Conclusion}
This paper presented a baseline-referenced feature-extraction method for single-channel EEG. Spectral and cross-band features are estimated from short EEG windows, aggregated into temporal blocks, and transformed into task-evoked deltas relative to an eyes-closed baseline. The primary output is a feature-delta vector whose direction and magnitude are characterized using percent change, standardized change, effect size, and statistical reliability. This formulation emphasizes within-participant task modulation while preserving explicit uncertainty and the limitations of single-channel measurement.

\begin{thebibliography}{99}

\bibitem{thomson1982}
D. J. Thomson, ``Spectrum estimation and harmonic analysis,'' \emph{Proceedings of the IEEE}, vol. 70, no. 9, pp. 1055--1096, 1982.

\bibitem{percival1993}
D. B. Percival and A. T. Walden, \emph{Spectral Analysis for Physical Applications}. Cambridge, U.K.: Cambridge Univ. Press, 1993.

\bibitem{whitham2007}
K. Whitham \emph{et al.}, ``Scalp electrical recording during paralysis: quantitative evidence that EEG frequencies above 20 Hz are contaminated by EMG,'' \emph{Clinical Neurophysiology}, vol. 118, pp. 1877--1888, 2007.

\bibitem{goncharova2003}
I. I. Goncharova, D. J. McFarland, T. M. Vaughan, and J. R. Wolpaw, ``EMG contamination of EEG: spectral and topographical characteristics,'' \emph{Clinical Neurophysiology}, vol. 114, pp. 1580--1593, 2003.

\bibitem{welch1947}
B. L. Welch, ``The generalization of Student's problem when several different population variances are involved,'' \emph{Biometrika}, vol. 34, no. 1/2, pp. 28--35, 1947.

\bibitem{cohen1988}
J. Cohen, \emph{Statistical Power Analysis for the Behavioral Sciences}, 2nd ed. Hillsdale, NJ, USA: Lawrence Erlbaum, 1988.

\bibitem{benjamini1995}
Y. Benjamini and Y. Hochberg, ``Controlling the false discovery rate: a practical and powerful approach to multiple testing,'' \emph{Journal of the Royal Statistical Society: Series B}, vol. 57, no. 1, pp. 289--300, 1995.

\bibitem{fisher1925}
R. A. Fisher, \emph{Statistical Methods for Research Workers}. Edinburgh, U.K.: Oliver and Boyd, 1925.

\bibitem{kost2002}
J. T. Kost and M. P. McDermott, ``Combining dependent $p$-values,'' \emph{Journal of Applied Statistics}, vol. 29, no. 1, pp. 225--240, 2002.

\bibitem{klimesch1999}
W. Klimesch, ``EEG alpha and theta oscillations reflect cognitive and memory performance: a review and analysis,'' \emph{Brain Research Reviews}, vol. 29, pp. 169--195, 1999.

\bibitem{kruskal1952}
W. H. Kruskal and W. A. Wallis, ``Use of ranks in one-criterion variance analysis,'' \emph{Journal of the American Statistical Association}, vol. 47, no. 260, pp. 583--621, 1952.

\bibitem{holm1979}
S. Holm, ``A simple sequentially rejective multiple test procedure,'' \emph{Scandinavian Journal of Statistics}, vol. 6, no. 2, pp. 65--70, 1979.

\end{thebibliography}

\end{document}
